\section*{第 \ref{ch:b1:fluid-envelopes-structure} 章——大气与海洋的结构}
% values: calc/ch26.py

\begin{solution}{exo:b1:fluid-envelopes-structure:1}
(a)~$H = R^*T/(M_0 g_0) = 8314.32 \times 250/(28.9644 \times 9.80665) =
\qty{7.32}{km}$。(b)~$e^{10/7.32} = 3.92$：压强约降至四分之一。
(c)~$H$ 与温度成正比：在 \qty{216.65}{K} 时为
\qty{6.34}{km}，而在 \qty{288.15}{K} 时为 \qty{8.43}{km}。冷空气在
同一压强下更稠密，因此压强随高度下降得更快。
\end{solution}

\begin{solution}{exo:b1:fluid-envelopes-structure:2}
取 $\Gamma_d = \qty{9.76}{K/km}$：\qtyrange{0}{0.3}{km} 内，$11.0 > 9.76$，
不稳定；\qtyrange{0.3}{1.2}{km} 内，$9.7$，在数据
精度范围内为中性；\qtyrange{1.2}{1.5}{km} 内为\omterm{def:b1:fluid-envelopes-structure:stability}{逆温}；\qtyrange{1.5}{8}{km} 内为稳定；
高于 \qty{8}{km} 时为等温、稳定。\omterm{def:b1:fluid-envelopes-structure:stability}{逆温}层最稳定，因为
$\Gamma_d - \Gamma = 13.8$ 在该处最大。
\end{solution}

\begin{solution}{exo:b1:fluid-envelopes-structure:3}
(a)~$\sigma = 1026.95 - 1000 = \qty{26.95}{kg.m^{-3}}$。(b)~较轻的 B 位于
A 之上。(c)~所有海水密度都以相同的数字 10 开头；$\sigma$ 只保留
在不同水体之间变化的那几位数字。
\end{solution}

\begin{solution}{exo:b1:fluid-envelopes-structure:4}
(a)~$S = 0.25 \times 35.40 + 0.75 \times 34.30 = 34.575$；$\theta = 0.25 \times
15.0 + 0.75 \times 4.0 = \qty{6.75}{\celsius}$。(b)~在线段 PQ 上，自 Q 向 P
方向的四分之一处（见\cref{prop:b1:fluid-envelopes-structure:mixing}）。
(c)~P、Q、R 的混合物充满三角形 PQR。令
$(34.60, 8.0) = a\,\mathrm P + b\,\mathrm Q + c\,\mathrm R$，且 $a + b + c = 1$
解得 $a = 0.345$，$b = 0.788$，$c = -0.133$：R 的比例为负，因此该
点位于三角形之外，三者的任何混合都无法得到它。
\end{solution}

\begin{solution}{exo:b1:fluid-envelopes-structure:5}
(a)~$\theta = 253.15 \times 2^{2/7} = \qty{308.6}{K}$。(b)~在 \qty{1000}{hPa} 处，
$\theta = T = \qty{298.15}{K}$。\omterm{def:b1:fluid-envelopes-structure:potential-temperature}{位温}随高度增加，由
298.2 增至 \qty{308.6}{K}：对于干空气，该层是稳定的。(c)~它将达到
其\omterm{def:b1:fluid-envelopes-structure:potential-temperature}{位温} \qty{308.6}{K}，即 \qty{35.4}{\celsius}，比
该处原有空气高出十度。
\end{solution}

\begin{solution}{exo:b1:fluid-envelopes-structure:6}
(a)~$T = 288.15 - 6.5 \times 5 = \qty{255.65}{K}$，且 $p = \num{101325} \times
(255.65/288.15)^{5.256} = \qty{54020}{Pa}$，即 \qty{540.2}{hPa}。(b)~$H =
\qty{8.43}{km}$，且 $p = \num{101325}\,e^{-5/8.43} = \qty{56010}{Pa}$。实际
（标准）压强更低：该处空气比 \qty{288.15}{K} 更冷，因而
更稠密，压强随高度下降得更快。(c)~$p = \num{22632}\,
e^{-9/6.34} = \qty{5475}{Pa}$，即为
\cref{tab:b1:fluid-envelopes-structure:layers}中的数值。
\end{solution}

\begin{solution}{exo:b1:fluid-envelopes-structure:7}
(a)~$\beta\,\Delta S = \alpha\,\Delta T$，得 $\Delta S = 1.67/7.61 = 0.22$。
(b)~$\Delta\rho = 1026.95 \times (-1.67\times10^{-4} \times 2 + 7.61\times10^{-4}
\times 0.3) = \qty{-0.11}{kg.m^{-3}}$：水变轻。
(c)~$\Delta\rho = 1028.11 \times (-0.53\times10^{-4} \times 2 + 7.85\times10^{-4}
\times 0.3) = \qty{+0.13}{kg.m^{-3}}$：同样的变化此时使水变重，因为
在冰点附近升温几乎不改变水的密度。
\end{solution}

\begin{solution}{exo:b1:fluid-envelopes-structure:8}
(a)~$N^2 = (9.81/260) \times (9.76 - 4)\times10^{-3} = \qty{2.17e-4}{s^{-2}}$，
$N = \qty{0.0147}{s^{-1}}$，周期 $2\pi/N = \qty{426}{s}$，约七
分钟。(b)~$N^2 = (9.81/270) \times 14.76\times10^{-3} = \qty{5.36e-4}{s^{-2}}$，
$N = \qty{0.0232}{s^{-1}}$，周期 \qty{271}{s}。(c)~$N$ 趋于零，周期
无限增长：在\omterm{def:b1:fluid-envelopes-structure:stability}{中性层结}中，回复力消失。
\end{solution}

\begin{solution}{exo:b1:fluid-envelopes-structure:9}
(a)~它更稠密，$\sigma_\theta = 27.20$，相比 26.30：AAIW 更淡，
但温度低得多。(b)~$\Delta\sigma = 1026.95 \times [-1.67\times10^{-4} \times
(4.31 - 15.43) + 7.61\times10^{-4} \times (34.30 - 35.54)] = 1.907 - 0.969 =
0.94$，与精确值 $0.90$ 相比。(c)~降温使
密度增加 \qty{1.91}{kg.m^{-3}}，淡化使密度减少 \qty{0.97}{kg.m^{-3}}。
\end{solution}

\begin{solution}{exo:b1:fluid-envelopes-structure:10}
(a)~由\cref{thm:b1:fluid-envelopes-structure:stability}，$(g/\theta)\,\mathrm
d\theta/\mathrm dz = (g/T)(\Gamma_d - \Gamma)$；$\mathrm d\theta/\mathrm dz = 0$
给出 $\Gamma = \Gamma_d$。(b)~$T = T_s - \Gamma_d z$，为梯度恒定的层，
其梯度 $L = \Gamma_d$：由\cref{prop:b1:fluid-envelopes-structure:layer-pressure}，$p =
p_s (T/T_s)^{gM/(R\Gamma_d)}$，且 $gM/(R\Gamma_d) = gMc_p/(Rg) = c_pM/R$
（对干空气，$= 3.5$）。(c)~$p$ 在 $T = 0$ 处为零，位于 $z = T_s/\Gamma_d =
300/9.76 = \qty{30.7}{km}$。\omterm{def:b1:fluid-envelopes-structure:potential-temperature}{位温}均匀的真实层至多
仅数千米厚：它们上方被稳定空气覆盖，正如对流层
被平流层覆盖。
\end{solution}

\begin{solution}{exo:b1:fluid-envelopes-structure:11}
(a)~$S = (33.787 + 36.050)/2 = 34.919$，$\theta = \qty{8.0}{\celsius}$。
(b)~对线性状态方程，$\sigma$ 为 $S$ 与
$\theta$ 的仿射函数，因此混合物的 $\sigma$ 为两者的平均值，即 27.00。
(c)~$\alpha$ 在 \qty{2}{\celsius} 时为 \qty{0.75e-4}{K^{-1}}，且
在 \qty{14}{\celsius} 时为 \qty{2.07e-4}{K^{-1}}，而 $\beta$ 变化很小
（分别为 \num{7.80e-4} 与 \num{7.53e-4}）。密度随
温度升高而下降得越来越快：它是温度的凹函数。沿混合线，
混合物处于平均温度，此处密度位于
连接两\omterm{def:b1:rock-forming-minerals:solid-solution}{端元}的弦之上，即高于 27.00。(d)~超出部分来自
曲率所致，是二阶项：它随两\omterm{def:b1:rock-forming-minerals:solid-solution}{端元}的温度
差的平方增长。
\end{solution}

\begin{solution}{exo:b1:fluid-envelopes-structure:12}
(a)~压强差为 $5604 - 5500 = \qty{104}{dbar}$：\omterm{def:b1:fluid-envelopes-structure:ts-diagram}{水团}降温
$0.129 \times 0.104 = \qty{0.013}{\celsius}$，降至 \qty{0.718}{\celsius}。
(b)~周围水体在 \qty{5400}{m} 处为 \qty{0.730}{\celsius}，更暖：该\omterm{def:b1:fluid-envelopes-structure:ts-diagram}{水团}
盐度几乎相同，比周围水更稠密，因而下沉
回原处。该层是稳定的，尽管实测温度随
深度增加，其\omterm{def:b1:fluid-envelopes-structure:potential-temperature}{位温}却随深度减小。(c)~$\theta \approx
0.7313 - 0.086 \times 5.604 = \qty{0.249}{\celsius}$，与精确值相差不超过
千分之一度。
\end{solution}

\begin{solution}{pb:b1:fluid-envelopes-structure:1}
\begin{enumerate}
  \item 温度在 \qty{50}{m} 深度以内几乎均匀，
        在 \qty{300}{m} 处迅速下降，较慢地降至 \qty{1000}{m} 处，以下几乎均匀。
        盐度自表层的 35.80 减小，
        至 \qty{800}{m} 处的最小值 34.32，然后增大到 \qty{3000}{m} 处的 34.90。
  \item $22.40 - 0.5 = \qty{21.90}{\celsius}$ 位于 \qty{50}{m}
        （22.10）与 \qty{75}{m}（19.20）之间：$50 + 25 \times 0.20/2.90 = \qty{52}{m}$。
  \item $(22.10 - 13.00)/250 = \qty{0.0364}{\celsius/m}$，相当于每
        \qty{100}{m} 降温 \qty{3.6}{\celsius}。\omterm{def:b1:fluid-envelopes-structure:ocean-layers}{温跃层}自\omterm{def:b1:fluid-envelopes-structure:ocean-layers}{混合层}底部
        （约 \qty{50}{m} 处）延伸至约 \qty{800}{m}；在
        \qty{50}{m} 与 \qty{75}{m} 之间最为陡峭（每 \qty{100}{m} 降温 \qty{11.6}{\celsius}）。
  \item $\sigma = 1026.95 \times [1 - 1.67\times10^{-4} \times 12.4 +
        7.61\times10^{-4} \times 0.80] - 1000 = 25.45$，偏高 \qty{0.74}{kg.m^{-3}}。
        表层水远离参考状态，且 $\alpha$
        在 \qty{22}{\celsius} 时远大于在 \qty{10}{\celsius} 时，
        因此线性形式低估了温度升高造成的变轻。
  \item \qty{800}{m} 处的盐度最小值 34.32（温度 \qty{4.8}{\celsius}）
        属于南极中层水，形成于南大洋多雨的
        西风带之下。\qty{3000}{m} 处的最大值 34.90（温度
        \qty{2.5}{\celsius}）属于北大西洋深层水，形成于
        格陵兰海与拉布拉多海。
  \item $(2.50 - 2.26)/3.042 = \qty{0.079}{\celsius}$，每 \qty{1000}{dbar}。
  \item $\sigma_\theta$ 随深度增大，该水柱是稳定的。
        \omterm{def:b1:fluid-envelopes-structure:ocean-layers}{密度跃层}位于 $\sigma_\theta$ 增大最快之处，即上部几
        百米之内，与\omterm{def:b1:fluid-envelopes-structure:ocean-layers}{温跃层}对应：在
        \qty{200}{m} 与 \qty{800}{m} 之间为 \num{0.82}，而
        其下方的 \qty{2200}{m} 内仅为 \num{0.70}。
  \item $S = (35.45 + 34.32)/2 = 34.885$; $\theta = (14.90 + 4.80)/2 =
        \qty{9.85}{\celsius}$.
  \item C：26.461；I：27.310；混合物：26.886，恰为两者的平均值
        （26.886），因为线性形式是仿射的。
  \item 两个母体的平均值为 $(26.342 + 27.159)/2 = 26.751$；混合物的密度
        为 26.888，比其平均值大 \qty{0.14}{kg.m^{-3}}。这就是混合增密效应。
  \item 按线性形式，该水柱在 \qty{300}{m} 处为 26.61，
        在 \qty{500}{m} 处为 26.88，在 \qty{800}{m} 处为 27.31：混合物的密度接近
        \qty{500}{m} 处。该处的水 $(34.82, \qty{9.6}{\celsius})$ 非常
        接近混合物 $(34.885, \qty{9.85}{\celsius})$：中央水与 AAIW 之间的水柱
        本身就近似位于
        两者之间的混合线上。
  \item $S = 0.25 \times 35.45 + 0.75 \times 34.32 = 34.60$, $\theta = 0.25
        \times 14.90 + 0.75 \times 4.80 = \qty{7.33}{\celsius}$。对
        \qty{500}{m} 与 \qty{800}{m} 之间的水柱做插值，得到
        $(34.57, \qty{7.2}{\celsius})$，在 \qty{650}{m} 附近，几乎正好落在混合线上。
  \item 以 \si{K/km} 表示：A--B 10.0；B--C 9.7；C--D $-5.5$；D--E 6.29；E--F 6.50；
        F--G 6.57; G--H 0.33; H--I $-0.50$; I--J $-0.83$.
  \item A--B 不稳定；B--C 中性；C--D 为\omterm{def:b1:fluid-envelopes-structure:stability}{逆温}；D--E、E--F、F--G 稳定；
        G--H 近似等温，稳定；H--I 与 I--J 为\omterm{def:b1:fluid-envelopes-structure:stability}{逆温}。对流层顶
        位于 G 处，高度 \qty{12.5}{km}，温度为 \qty{-59.0}{\celsius}，在此
        温度不再下降。
  \item 平均温度为 \qty{214.65}{K}：$H = 287 \times 214.65/9.81 =
        \qty{6.28}{km}$，且 $147\,e^{-4/6.28} = \qty{77.7}{hPa}$，与
        实测的 78 相符。
  \item A 处的空气温度为 \qty{22.0}{\celsius}，低于海水温度：热量由
        海水流向空气，空气从下方受热。A--B 层，其
        递减率大于 $\Gamma_d$，正是这一过程的结果。
  \item 在 B 处：$22.0 - 9.76 \times 0.5 = \qty{17.1}{\celsius}$，对比 17.0；在
        C 处：\qty{7.4}{\celsius}，对比 7.3；在 D 处：\qty{3.5}{\celsius}，对比
        9.5。直到 C 处，\omterm{def:b1:fluid-envelopes-structure:parcel}{气块}都略比空气暖，
        而在 D 处则冷得多。
  \item 在\omterm{def:b1:fluid-envelopes-structure:stability}{逆温}层底部，即 \qty{1.5}{km} 附近：进入\omterm{def:b1:fluid-envelopes-structure:stability}{逆温}层时，
        \omterm{def:b1:fluid-envelopes-structure:parcel}{气块}变得比周围空气更冷、因而更稠密，
        于是停止上升。
  \item 不稳定表层释放的\omterm{def:b1:fluid-envelopes-structure:parcel}{气块}向上穿过该层并使其混合，
        直至其\omterm{def:b1:fluid-envelopes-structure:potential-temperature}{位温}均匀，使递减率
        为 $\Gamma_d$；上方的\omterm{def:b1:fluid-envelopes-structure:stability}{逆温}层阻止混合继续进行。
  \item 以 K 表示：A 294.15；B 293.99；C 294.08；D 300.51；E 312.65；F 329.29；G
        346.28; H 369.50; I 446.99; J 600.91.
  \item $\theta$ 自 A 到 B 减小（不稳定），自 B 到
        C 近似不变（中性），再往上增大（稳定），自 C 到 D 以及
        G 以上增大最快：与练习~14 的分类相同。
  \item E--F: $(329.29 - 312.65)/4 = \qty{4.2}{K/km}$; H--I: $(446.99 -
        369.50)/4 = \qty{19.4}{K/km}$，约为前者的五倍。
  \item 取平均\omterm{def:b1:fluid-envelopes-structure:potential-temperature}{位温} $\theta = \qty{408.2}{K}$：$N^2 = (9.81/408.2) \times
        0.01937 = \qty{4.66e-4}{s^{-2}}$，$N = \qty{0.0216}{s^{-1}}$，周期
        为 \qty{291}{s}，约五分钟。
  \item $N^2 = (9.81/1027) \times (26.35 - 24.79)/150 = \qty{9.9e-5}{s^{-2}}$,
        $N = \qty{0.0100}{s^{-1}}$，周期为 \qty{630}{s}。平流层下部
        更为稳定，其 $N^2$ 约为
        \omterm{def:b1:fluid-envelopes-structure:ocean-layers}{温跃层}中的五倍。
  \item 在 G 处，$\theta = 214.15 \times (1000/186)^{2/7} = \qty{346}{K}$：
        对流层顶处的\omterm{def:b1:fluid-envelopes-structure:potential-temperature}{位温}约为 \qty{346}{K}。
\end{enumerate}
\end{solution}
